% ==========================================================================
%  Decoreba V/F — P2 2026 (guia de memorização de emergência)
%  Macroeconomia I (PPGEco-UFES, 2026)
%  Para gravar: gabaritos literais, armadilhas, gatilhos, pares, autores,
%  números e drill de 54 afirmativas relâmpago.
%  Fontes: P2 2021, P2 2024 (chave oficial), Lista 3 2026 (chave oficial),
%          slides das Aulas 10-13.
% ==========================================================================
\documentclass[12pt, a4paper]{article}

\usepackage[utf8]{inputenc}
\usepackage[T1]{fontenc}
\usepackage{lmodern}
\usepackage[brazil]{babel}
\usepackage{microtype}

\usepackage[top=2.2cm, bottom=2.2cm, left=2.4cm, right=2.4cm, headheight=15pt]{geometry}

\usepackage{amsmath, amssymb}
\usepackage{mathtools}

\usepackage{booktabs}
\usepackage{array}
\usepackage{xcolor}
\usepackage{hyperref}
\usepackage{enumitem}
\usepackage{tcolorbox}
\tcbuselibrary{breakable}
\usepackage{fancyhdr}
\usepackage{titlesec}
\usepackage{tikz}

% --------------------------------------------------------------------------
% Cores
% --------------------------------------------------------------------------
\definecolor{cyanrbc}{HTML}{3b9dbd}
\definecolor{orangerbc}{HTML}{d97b39}
\definecolor{grayrbc}{HTML}{6b7280}
\definecolor{greenbox}{HTML}{166534}
\definecolor{greenbg}{HTML}{dcfce7}
\definecolor{boxbg}{HTML}{f5f0e8}
\definecolor{boxborder}{HTML}{3b9dbd}
\definecolor{provabg}{HTML}{fff7ed}
\definecolor{provaborder}{HTML}{d97b39}
\definecolor{correctbg}{HTML}{dcfce7}
\definecolor{correctborder}{HTML}{166534}
\definecolor{redf}{HTML}{b91c1c}
\definecolor{redbg}{HTML}{fee2e2}

% --------------------------------------------------------------------------
% Caixas
% --------------------------------------------------------------------------
\newtcolorbox{resultbox}[1][]{standard, colback=boxbg, colframe=boxborder,
  boxrule=1pt, arc=3pt, left=6pt, right=6pt, top=4pt, bottom=4pt,
  fonttitle=\bfseries\small, breakable, #1}
\newtcolorbox{provabox}[1][]{standard, colback=provabg, colframe=provaborder,
  boxrule=1.2pt, arc=3pt, left=6pt, right=6pt, top=4pt, bottom=4pt,
  fonttitle=\bfseries\small, breakable, #1}
\newtcolorbox{correctbox}[1][]{standard, colback=correctbg, colframe=correctborder,
  boxrule=1pt, arc=3pt, left=6pt, right=6pt, top=4pt, bottom=4pt,
  fonttitle=\bfseries\small\color{correctborder}, breakable, #1}
\newtcolorbox{falsebox}[1][]{standard, colback=redbg!50, colframe=redf!80,
  boxrule=1pt, arc=3pt, left=6pt, right=6pt, top=4pt, bottom=4pt,
  fonttitle=\bfseries\small\color{redf}, breakable, #1}

% --------------------------------------------------------------------------
% Comandos
% --------------------------------------------------------------------------
\newcommand{\E}{\mathbb{E}}
\newcommand{\dif}{\mathrm{d}}
\newcommand{\VAR}{\mathrm{VAR}}
\newcommand{\Vok}{\textbf{\textcolor{greenbox}{V}}}
\newcommand{\Fno}{\textbf{\textcolor{redf}{F}}}
\newcommand{\mata}[1]{\textcolor{redf}{\underline{\textbf{#1}}}}

\setlength{\parskip}{0.45\baselineskip}
\setlength{\parindent}{0pt}

\titleformat{\section}{\large\bfseries\color{cyanrbc}}{\thesection.}{0.5em}{}[\vspace{-0.3em}\rule{\linewidth}{0.4pt}]

\pagestyle{fancy}
\fancyhf{}
\lhead{\footnotesize\textcolor{grayrbc}{Decoreba V/F --- P2 2026 --- Macroeconomia I (PPGEco-UFES)}}
\rhead{\footnotesize\textcolor{grayrbc}{ler $\cdot$ falar alto $\cdot$ repetir}}
\cfoot{\small\thepage}
\renewcommand{\headrulewidth}{0.4pt}

% ==========================================================================
\begin{document}
% ==========================================================================

% --------------------------------------------------------------------------
% Capa compacta
% --------------------------------------------------------------------------
\begin{center}
\vspace*{0.2cm}
\rule{\linewidth}{1.2pt}\\[0.6em]
{\LARGE\bfseries\color{cyanrbc} Decoreba V/F --- Prova 2}\\[0.3em]
{\large Macroeconomia I --- PPGEco-UFES 2026}\\[0.3em]
{\normalsize\color{grayrbc} Guia de memorização de emergência
(caps.\ 8 $\cdot$ 9 $\cdot$ 11 $\cdot$ 12)}\\[0.5em]
\rule{\linewidth}{1.2pt}
\end{center}

\begin{provabox}[title={Modo de usar (sem tempo para teoria)}]
\begin{enumerate}[leftmargin=1.6em, itemsep=0.1em, topsep=0.2em]
\item Grave as \textbf{7 frases de ouro} (§1) --- elas \emph{foram} gabarito
oficial e tendem a voltar;
\item Grave os \textbf{falsos clássicos} (§2) pela \mata{palavra que mata};
\item Decore a tabela de \textbf{gatilhos} (§3) e os \textbf{pares} (§4);
\item Faça o \textbf{drill de 54} (§8) marcando V/F com a mão --- 3 rodadas;
\item 10 minutos antes da prova: releia só §1, §3 e o quadro final (§9).
\end{enumerate}
\end{provabox}

% ------------------------------------------------------------------
\section{As 7 frases de ouro (foram gabarito --- decore como VERDADEIRAS)}
% ------------------------------------------------------------------

\begin{correctbox}[title={Se aparecer parecido com isto, é V}]
\begin{enumerate}[leftmargin=1.6em, itemsep=0.35em]
\item \textbf{[Consumo]} ``A rejeição \emph{plena} da hipótese do passeio
aleatório requer encontrar valores convergentes com $\lambda = 1$.''
\hfill{\footnotesize (P2 2021 e P2 2024, Q1)}
\item \textbf{[Consumo/Fiscal]} ``A restrição orçamentária das famílias, com
presença do governo, \emph{não} sofre influência das proporções de $T$ e $D$.''
\hfill{\footnotesize (P2 2024, Q2)}
\item \textbf{[Investimento]} ``$K$ da indústria muito baixo $+$ $q$ também
muito reduzido $\Rightarrow$ \emph{queda adicional} do $q$.''
[$\dot q = rq - \pi(K) < 0$]
\hfill{\footnotesize (P2 2021 Q2 e P2 2024 Q3)}
\item \textbf{[Investimento]} ``Mais importante que o juro real \emph{de hoje}
é a expectativa de $r$ \emph{para o longo prazo}, no que respeita ao $q$.''
\hfill{\footnotesize (P2 2024, Q4)}
\item \textbf{[Moeda]} ``Preços flexíveis $+$ redução permanente do
crescimento da moeda $\Rightarrow$ queda \emph{descontínua} de $P$, de
$\pi^e$ e de $i$, com \emph{aumento} de $M/P$ no longo prazo.''
\hfill{\footnotesize (P2 2021, Q3)}
\item \textbf{[Fiscal]} ``Satisfeita a restrição do governo, substituir
financiamento via tributos por mais dívida \emph{não} impacta o consumidor.''
\hfill{\footnotesize (P2 2024, Q5)}
\item \textbf{[Fiscal]} ``Os superávits primários \emph{corrente e futuros
esperados} precisam cobrir a dívida pública atual em termos reais.''
\hfill{\footnotesize (P2 2024, Q6)}
\end{enumerate}
\end{correctbox}

% ------------------------------------------------------------------
\section{Falsos clássicos (decore a palavra que mata)}
% ------------------------------------------------------------------

\begin{falsebox}[title={Se aparecer parecido com isto, é F --- o erro está sublinhado}]
\textbf{Consumo}
\begin{enumerate}[leftmargin=1.6em, itemsep=0.15em, topsep=0.2em]
\item ``Em amostras \mata{longas}, a variância da renda \mata{transitória}
domina'' --- é o contrário: amostra longa $\Rightarrow$ domina a
\emph{permanente} (e $\hat b \to 1$).
\item ``Correlação juros--consumo implica o \mata{abandono} da hipótese da
renda permanente'' --- a equação de Euler é \emph{extensão} dela, não
abandono.
\item ``Poupança precaucional ajuda a entender a \mata{elevada} poupança das
famílias'' --- o fato estilizado é poupança \emph{baixa} (puzzle; Carroll).
\item ``Variação do consumo \mata{inversamente} relacionada à taxa real de
juros'' --- é \emph{positiva} (Euler).
\item ``Há \mata{grande} impacto dos juros reais sobre o consumo'' --- o
impacto empírico é \emph{fraco} ($\theta$ alto).
\item ``Maior aversão ao risco $\Rightarrow$ \mata{maior} eficácia dos juros
sobre o consumo'' --- maior $\theta$ $\Rightarrow$ \emph{menor} sensibilidade
($1/\theta$).
\end{enumerate}
\textbf{Investimento}
\begin{enumerate}[leftmargin=1.6em, itemsep=0.15em, topsep=0.2em]
\setcounter{enumi}{6}
\item ``A firma investe até o valor do capital \mata{superar} o custo de
aquisição'' --- no ótimo \emph{iguala}: $1 + C'(I) = q$.
\item ``Abaixo de $\dot K=0$ e à direita de $\dot q=0$ é região de
\mata{divergência} crescente'' --- ali passa o braço \emph{convergente} da
trajetória de sela.
\item ``Redução transitória do produto tem efeito \mata{nulo} sobre o
investimento'' --- efeito \emph{menor}, nunca nulo.
\item ``Valor da empresa \mata{positivamente} relacionado ao $K$ da
indústria'' / ``$\dot q = 0$ \mata{positivamente} inclinada'' --- $\pi'(K)<0$:
relação e inclinação \emph{negativas}.
\item ``Inflação esperada acima da meta $\Rightarrow$ \mata{maior} $q$'' ---
Banco Central sobe o juro real ($\phi_\pi>1$) $\Rightarrow$ $q$ \emph{cai}.
\end{enumerate}
\textbf{Moeda}
\begin{enumerate}[leftmargin=1.6em, itemsep=0.15em, topsep=0.2em]
\setcounter{enumi}{11}
\item ``Preços sobem com variações positivas na moeda, nos juros e no
\mata{produto}'' --- produto entra \emph{negativo} ($L_Y>0$ eleva demanda por
moeda e derruba $P$).
\item ``Arbitragem: $i^n$ igual a \mata{$n$ vezes $i^1$}'' --- $i^n$ é a
\emph{média} das taxas curtas esperadas $+$ prêmio a termo.
\item ``Efeito liquidez: expansão monetária \mata{eleva} as taxas curtas'' ---
\emph{reduz} as curtas (e tende a subir as \emph{longas}).
\item ``Com $\phi_\pi < 1$ a inflação converge, só que \mata{devagar}'' ---
não converge: fica \emph{persistente} e as expectativas \emph{desancoram}.
\end{enumerate}
\textbf{Fiscal}
\begin{enumerate}[leftmargin=1.6em, itemsep=0.15em, topsep=0.2em]
\setcounter{enumi}{15}
\item ``Solvência avaliada pelos resultados fiscais \mata{correntes}'' ---
vale o fluxo \emph{projetado} (expectativas, anúncios, credibilidade).
\item ``B com pior saldo corrente é \mata{necessariamente} mais arriscado''
--- expectativas de fluxos futuros podem inverter a comparação.
\item ``Avaliação fiscal \mata{não sofre qualquer influência} da inflação''
--- sofre; o item generaliza demais.
\item ``Redução de alíquotas $\Rightarrow$ \mata{aumento do consumo} e
estímulo real'' --- equivalência ricardiana: a renda extra é \emph{poupada}.
\item ``Gastos públicos \mata{não} têm impacto direto na restrição das
famílias'' --- têm: após substituir, $-\int e^{-R}G$ entra na restrição.
\item ``São os \mata{tributos} que afetam a acumulação de capital'' --- são os
\emph{gastos} ($I = Y - C - G$).
\item ``Alesina--Perotti: ajustes via \mata{tributos} são os duradouros e
expansionistas'' --- são os via \emph{cortes de gastos} (emprego público e
transferências).
\item ``Calvo/Cole--Kehoe: dívida alta $\Rightarrow$ retorno requerido
\mata{menor}'' --- \emph{maior} (e default mais provável; crise
autorrealizável).
\end{enumerate}
\end{falsebox}

% ------------------------------------------------------------------
\section{Gatilhos de marcação (regras de bolso)}
% ------------------------------------------------------------------

\begin{resultbox}[title={Palavra no item $\Rightarrow$ chute educado (e a exceção)}]
\renewcommand{\arraystretch}{1.25}
\begin{tabular}{p{4.6cm} c p{8.6cm}}
\toprule
\textbf{Se o item diz\ldots} & \textbf{Chute} & \textbf{Por quê / exceção} \\
\midrule
``supera'' / ``excede'' (no ótimo) & \Fno & condição de ótimo é
\emph{igualdade} ($1{+}C'(I){=}q$) \\
``efeito nulo'' (choque transitório) & \Fno & efeito é \emph{menor}, nunca
nulo \\
``necessariamente'' & \Fno & generalização; expectativas podem inverter \\
``abandono da teoria'' & \Fno & extensões não são abandono \\
``elevada poupança das famílias'' & \Fno & fato estilizado: poupam
\emph{pouco} \\
``grande impacto dos juros no consumo'' & \Fno & impacto empírico é
\emph{fraco} \\
``não sofre qualquer influência'' & \Fno & \textbf{exceto} se for sobre $T$ e
$D$ na restrição das famílias (ricardiana) --- aí é \Vok \\
``backward-looking'' (Clarida) & \Fno & a regra é \emph{forward-looking}
(expectativas) \\
``positivamente inclinada'' ($\dot q=0$) & \Fno & $\pi'(K)<0$: inclinação
\emph{negativa} \\
``$i^n = n\,i^1$'' & \Fno & é \emph{média} $+$ prêmio a termo \\
``média das taxas curtas esperadas'' & \Vok & estrutura a termo correta \\
``iguala'' (custo vs.\ valor no ótimo) & \Vok & é a condição de otimalidade \\
``queda adicional do $q$'' ($K$ e $q$ baixos) & \Vok & $\dot q = rq -
\pi(K) < 0$ \\
``juros de longo prazo importam para o $q$'' & \Vok & $r$ é a taxa de
desconto dos fluxos \\
\bottomrule
\end{tabular}
\end{resultbox}

% ------------------------------------------------------------------
\section{Pares que o professor troca (decore os dois lados)}
% ------------------------------------------------------------------

\begin{resultbox}[title={Conceito certo $\times$ troca falsa}]
\renewcommand{\arraystretch}{1.25}
\begin{tabular}{p{6.6cm} p{6.9cm}}
\toprule
\textbf{Certo} & \textbf{Troca que vira F} \\
\midrule
$\theta$ = aversão ao risco; EIS $= 1/\theta$ &
dizer que $\theta$ \emph{é} a EIS \\
$(\theta + 1)$ = \emph{prudência} relativa &
chamar $(\theta+1)$ de aversão ao risco \\
transitório: impacto $1/T$ no consumo &
dizer que transitório desloca 1-para-1 \\
permanente: impacto $\approx 1$ &
dizer que permanente tem impacto $1/T$ \\
amostra curta: $\hat b \ll 1$ (transitória pesa) &
inverter qual variância domina \\
amostra longa: $\hat b \to 1$ (permanente domina) &
--- \\
$\lambda \approx 0{,}5$, significativo = híbrido &
dizer que acharam $\lambda = 1$ \\
Hall: informação \emph{já incorporada} &
dizer que famílias \emph{ignoram} informação \\
Campbell--Mankiw: \emph{IV} (endogeneidade) &
dizer que foi \emph{MQO} \\
$r > \rho \Rightarrow$ consumo \emph{cresce} &
inverter para decrescer \\
$\dot K = 0$: reta em $q = 1$ &
dizer $q = 0$ ou ``cresce à taxa $r$'' \\
$\dot q = 0$: $q = \pi(K)/r$, \emph{decrescente} &
inclinação positiva \\
NO e SE do diagrama: \emph{sela} (converge) &
chamar de divergência \\
NE e SO: divergem &
--- \\
Fisher: $i = r + \pi^e$ (nominal carrega $\pi^e$) &
escrever $r = i + \pi^e$ \\
efeito liquidez: \emph{curtas caem} &
dizer que curtas sobem \\
longas \emph{sobem} (expectativas de curtas) &
atribuir ao prêmio a termo \\
Taylor: $\phi_\pi > 1 \Rightarrow$ \emph{real} sobe &
aplicar o ``mais que proporcional'' ao real \\
médio prazo: $i = r^n + \pi^*$ ($5{+}3{=}8\%$) &
escrever $i = r^n - \pi^*$ \\
solvência: fluxos \emph{esperados} a VP &
avaliar só o corrente \\
ricardiana: \emph{G} importa; $T$/$D$ somem &
dizer que tributos movem capital \\
Alesina--Perotti: corte de \emph{gastos} expande &
dizer que é o aumento de tributos \\
\bottomrule
\end{tabular}
\end{resultbox}

% ------------------------------------------------------------------
\section{Quem fez o quê (autores caem nos itens!)}
% ------------------------------------------------------------------

\begin{resultbox}[title={Nome $\to$ uma linha}]
\renewcommand{\arraystretch}{1.2}
\begin{tabular}{p{5.2cm} p{8.3cm}}
\toprule
Keynes (1936) & consumo $=$ função da renda disponível corrente \\
Modigliani--Brumberg (1954) & ciclo de vida (com Friedman: renda permanente) \\
Friedman (1957) & renda permanente; Keynes vira caso particular \\
Hall (1978) & passeio aleatório: $C_t = C_{t-1} + e_t$ \\
Hansen--Singleton (1983); Hall (1988b) & juros reais explicam pouco o
consumo \\
Campbell--Mankiw (1989) & teste IV; $\lambda \approx 0{,}5$; explicação
híbrida \\
Carroll (1992, 1997) & desconto alto anula o motivo precaucional \\
Zeldes (1989) & restrição de liquidez \emph{futura} já reduz consumo \\
Gross--Souleles (2002) & evidência de restrições de crédito \\
Tobin (1969) & $q$: valor de mercado do capital \\
Taylor (1993) & regra original; princípio $\phi_\pi > 1$ \\
Clarida--Galí--Gertler (2000) & regra \emph{forward-looking} \\
Feldstein (1997) & inflação distorce tributos (regressiva, alimentos) \\
Bénabou (1992); Tommasi (1994) & inflação reduz eficiência \\
Modigliani--Cohn (1979); Hall (1984) & agentes erram contas com inflação \\
Giavazzi--Pagano (1990) & Dinamarca/Irlanda: contração fiscal expansionista \\
Alesina--Perotti (1997) & composição do ajuste: gastos $>$ tributos \\
Bertola--Drazen (1993); Perotti (1999) & renda futura esperada $\to$ consumo
e investimento \\
Calvo (1988); Cole--Kehoe (2000) & crises de dívida autorrealizáveis \\
\bottomrule
\end{tabular}
\end{resultbox}

% ------------------------------------------------------------------
\section{Fórmulas faladas (leia em voz alta 3 vezes)}
% ------------------------------------------------------------------

\begin{resultbox}[title={Fórmula $+$ leitura $+$ o sinal que importa}]
\begin{itemize}[leftmargin=1.4em, itemsep=0.4em]
\item $C_t = \frac{1}{T}\bigl(A_0 + \sum_\tau Y_\tau\bigr)$ ---
``consumo é \emph{um sobre T} da riqueza mais a soma das rendas'':
transitório mexe $1/T$; permanente mexe $\approx 1$.
\item $\hat b = \frac{\VAR(Y^P)}{\VAR(Y^P) + \VAR(Y^T)}$ ---
``permanente em cima, tudo embaixo'': quem domina o numerador manda o
coeficiente para 1.
\item $C_t = C_{t-1} + e_t$ --- ``consumo de hoje é o de ontem mais
surpresa'': só \emph{informação nova} move.
\item $\Delta C_t = \lambda\,\Delta Y_t + (1-\lambda)e_t$ ---
``lambda mede os keynesianos'': deu $\approx 0{,}5$, significativo, por IV.
\item $\frac{C_{t+1}}{C_t} = \bigl(\frac{1+r}{1+\rho}\bigr)^{1/\theta}$ ---
``juro acima da pressa, consumo em rampa''; expoente $1/\theta$ = EIS
(sensibilidade).
\item $\E[g^c] \simeq \frac{1}{2}(\theta+1)\VAR(g^c)$ ---
``meio, prudência, variância'': mais incerteza $\Rightarrow$ mais poupança
$\Rightarrow$ consumo cresce \emph{depois}.
\item $1 + C'(I) = q$ --- ``um mais custo marginal \emph{iguala} q'' (nunca
``supera'').
\item $q(t) = \int_0^\infty e^{-rt}\pi(K(t))\,\dif t$ ---
``q é o VP dos lucros'': sobe com produto esperado, cai com juros (longos) e
com $K$ da indústria.
\item $\dot q = rq - \pi(K)$; $\dot q = 0 \Rightarrow q = \pi(K)/r$ ---
``custo de oportunidade menos lucro'': $K$ e $q$ baixos $\Rightarrow$
$\dot q < 0$ (q cai mais).
\item $P = M/L(i,Y)$ --- ``M e i incham P; Y murcha P''.
\item $i = r + \pi^e$ --- ``o \emph{nominal} carrega a inflação''.
\item $i_t^n = \frac{1}{n}\sum_{j=0}^{n-1}\E_t\,i^1_{t+j} + \theta_{nt}$ ---
``longa é média das curtas esperadas mais prêmio''.
\item $i_t = r^n + \E_t\pi_{t+n} + \phi_\pi(\E_t\pi_{t+n} - \pi^*) +
\phi_y(\text{hiato esperado})$ --- ``natural, inflação esperada, desvio vezes
phi-pi, hiato vezes phi-y''; $\phi_\pi > 1$ para o \emph{real} subir.
\item $\int_0^\infty e^{-R(t)}[T(t)-G(t)]\,\dif t \geq D(0)$ ---
``superávits a valor presente cobrem a dívida de hoje''.
\item $\int e^{-R}C \leq K(0) + \int e^{-R}W - \int e^{-R}G$ ---
``depois da conta, sobram K, salários e \emph{G}: T e D sumiram''
(equivalência ricardiana).
\end{itemize}
\end{resultbox}

% ------------------------------------------------------------------
\section{O diagrama de fases de bolso (uma figura para gravar)}
% ------------------------------------------------------------------

\begin{center}
\begin{tikzpicture}[scale=1.0, >=stealth]
% eixos
\draw[->] (0,0) -- (9.2,0) node[below] {$K$};
\draw[->] (0,0) -- (0,5.6) node[left] {$q$};
% locus Kdot = 0 (q = 1)
\draw[thick, cyanrbc] (0,2.2) -- (8.8,2.2) node[right, font=\small] {$\dot K = 0$ ($q=1$)};
% locus qdot = 0 (decrescente)
\draw[thick, orangerbc] (0.7,5.1) .. controls (3.4,3.0) and (5.4,2.5) .. (8.6,1.1)
  node[right, font=\small] {$\dot q = 0$};
% ponto E
\filldraw (4.55,2.2) circle (2pt) node[above right, font=\small] {$E$};
% trajetoria de sela
\draw[dashed, thick, greenbox] (1.3,4.4) .. controls (3.2,3.0) .. (4.55,2.2)
  .. controls (6.0,1.55) .. (8.2,1.0);
\node[greenbox, font=\footnotesize, align=center] at (2.15,4.7) {sela\\(converge)};
\node[greenbox, font=\footnotesize, align=center] at (7.6,0.55) {sela\\(converge)};
% setas de dinamica por quadrante
% NO: K up, q down
\draw[->, thick] (1.7,3.3) -- (2.3,3.3); \draw[->, thick] (1.7,3.3) -- (1.7,2.85);
\node[font=\footnotesize, align=left] at (1.95,3.75) {NO: $\dot K{>}0$, $\dot q{<}0$};
% NE: K up, q up (diverge)
\draw[->, thick] (6.7,3.3) -- (7.3,3.3); \draw[->, thick] (6.7,3.3) -- (6.7,3.75);
\node[font=\footnotesize, redf, align=left] at (7.15,4.15) {NE: diverge $\nearrow$};
% SO: K down, q down (diverge)
\draw[->, thick] (2.0,1.2) -- (1.4,1.2); \draw[->, thick] (2.0,1.2) -- (2.0,0.75);
\node[font=\footnotesize, redf, align=left] at (2.5,0.5) {SO: diverge $\swarrow$ \ (``$q$ cai mais'')};
% SE: K down, q up
\draw[->, thick] (6.9,1.55) -- (6.3,1.55); \draw[->, thick] (6.9,1.55) -- (6.9,2.0);
\node[font=\footnotesize, align=left] at (6.35,1.1) {SE: $\dot K{<}0$, $\dot q{>}0$};
\end{tikzpicture}
\end{center}

\begin{resultbox}[title={Mnemônico do mapa}]
\textbf{``NO--SE converge, NE--SO explode.''} Acima de $q=1$, $K$ anda para a
\emph{direita}; abaixo, para a \emph{esquerda}. À direita de $\dot q = 0$, $q$
sobe; à esquerda, desce. A questão da Lista (Q5) é o quadrante \textbf{NO}:
$q>1$ ($\dot K > 0$) e pessimismo ($\dot q<0$) $\Rightarrow$ desce pela sela
até $E$. A frase de ouro n\textsuperscript{o}~3 é o quadrante \textbf{SO}:
$K$ baixo $+$ $q$ baixo $\Rightarrow$ $q$ cai ainda mais.
\end{resultbox}

% ------------------------------------------------------------------
\section{Números e fatos secos (decorar mesmo)}
% ------------------------------------------------------------------

\begin{resultbox}[title={Cartela de números}]
\begin{itemize}[leftmargin=1.4em, itemsep=0.15em]
\item Campbell--Mankiw: dados \textbf{trimestrais}, EUA,
\textbf{1953--1986}, $\lambda \approx \textbf{0,5}$ (significativo,
$\ll 1$), método \textbf{IV}.
\item Correlação consumo $\times$ PIB no Brasil (slides):
\textbf{0,99} (1996T1--2025T1, SCN/IBGE).
\item Calibração da regra (linha BCB): $r^n = \textbf{5\%}$;
$\pi^* = \textbf{3\%}$; $n = \textbf{4 trimestres}$;
$\phi_\pi = \textbf{1,5}$; $\phi_y = \textbf{0,3}$;
médio prazo $\Rightarrow i = \textbf{8\%}$.
\item IRFs em aula: $\phi_\pi = 1{,}5$ derruba a inflação \emph{mais rápido}
que $\phi_\pi = 0{,}7$ (ao custo de queda maior do produto).
\item Brasil fiscal (slides, BCB): primário médio $\textbf{-0,34\%}$ do PIB e
nominal $\approx \textbf{-4,5\%}$ (06/1999--05/2025); dívida bruta:
$\approx$\textbf{55\%} (2006) $\to$ pico $\approx$\textbf{90\%} (2020) $\to$
$\approx$\textbf{75\%} (2025).
\item Contrações expansionistas: \textbf{Dinamarca e Irlanda, anos 1980}
(Giavazzi--Pagano, 1990).
\item Inflação $\times$ crescimento monetário: \textbf{97 países,
1980--2006}, relação clara e forte (Fig.\ 11.1).
\end{itemize}
\end{resultbox}

% ------------------------------------------------------------------
\section{Drill relâmpago: 54 afirmativas (marque V ou F; 3 rodadas)}
% ------------------------------------------------------------------

Cubra a coluna da direita, julgue, confira. Erro $=$ voltar ao parágrafo
correspondente nos §§1--7.

\subsection*{Bloco A --- Consumo}
{\footnotesize
\renewcommand{\arraystretch}{1.3}
\begin{tabular}{p{0.5cm} p{9.2cm} c p{3.6cm}}
\toprule
& \textbf{Afirmativa} & & \textbf{Porquê} \\
\midrule
1 & O consumo de cada período é fração $1/T$ dos recursos de vida inteira. &
\Vok & solução do Lagrangiano \\
2 & Um ganho transitório de renda desloca o consumo um-para-um. &
\Fno & impacto é $1/T$ \\
3 & Aumento permanente da renda eleva o consumo $\approx$ um-para-um. &
\Vok & soma $T$ parcelas $1/T$ \\
4 & Em amostras longas a propensão estimada aproxima-se de 1. &
\Vok & $\VAR(Y^P)$ domina \\
5 & Em amostras curtas domina a variância da renda permanente. &
\Fno & domina a transitória \\
6 & Hall: consumo imprevisível porque a informação já está toda no consumo
corrente. & \Vok & $C_t = \E_t[C_{t+1}]$ \\
7 & Campbell--Mankiw estimaram por MQO. & \Fno & IV (endogeneidade) \\
8 & $\lambda$ significativo e $\approx 0{,}5$ $\Rightarrow$ explicação
híbrida. & \Vok & metade regra de bolso \\
9 & Rejeição plena do passeio aleatório requer $\lambda = 1$. &
\Vok & frase de ouro 1 \\
10 & Correlação juros--consumo obriga abandonar a renda permanente. &
\Fno & Euler é extensão \\
11 & $r > \rho$ $\Rightarrow$ consumo crescente no tempo. &
\Vok & Euler \\
12 & Aversão maior $\Rightarrow$ consumo mais sensível aos juros. &
\Fno & sensibilidade $= 1/\theta$ \\
13 & Evidência: juros reais têm forte impacto no consumo. &
\Fno & impacto fraco ($\theta$ alto) \\
14 & $(\theta + 1)$ é o coeficiente de prudência relativa. &
\Vok & não é aversão \\
15 & Mais incerteza ($\VAR(g^c)\uparrow$) $\Rightarrow$ $\E[g^c]$ maior. &
\Vok & poupa hoje, cresce depois \\
16 & O motivo precaucional explica a elevada poupança observada. &
\Fno & poupam pouco (puzzle) \\
17 & Carroll: taxa de desconto alta anula o motivo precaucional. &
\Vok & concilia teoria e fato \\
18 & Restrição de liquidez só importa se estiver ativa hoje. &
\Fno & Zeldes: futura já reduz \\
\bottomrule
\end{tabular}
}

\subsection*{Bloco B --- Investimento}
{\footnotesize
\renewcommand{\arraystretch}{1.3}
\begin{tabular}{p{0.5cm} p{9.2cm} c p{3.6cm}}
\toprule
19 & No ótimo, $1 + C'(I) = q$. & \Vok & custo iguala valor \\
20 & A firma investe até o valor do capital superar o custo. &
\Fno & ``supera'' mata \\
21 & $\dot K = 0 \iff q = 1$. & \Vok & $f(1) = 0$ \\
22 & O locus $\dot q = 0$ é positivamente inclinado. &
\Fno & $\pi'(K) < 0$ \\
23 & O valor da empresa cresce com o $K$ da indústria. &
\Fno & $\pi'(K) < 0$ \\
24 & O $q$ cai quando os juros longos esperados sobem. &
\Vok & $r$ desconta os fluxos \\
25 & $K$ muito baixo $+$ $q$ muito baixo $\Rightarrow$ $q$ cai ainda mais. &
\Vok & frase de ouro 3 \\
26 & Abaixo de $\dot K{=}0$ e à direita de $\dot q{=}0$: divergência
crescente. & \Fno & braço da sela (SE) \\
27 & Acima de $\dot K{=}0$ e à esquerda de $\dot q{=}0$: converge a $E$ pela
sela. & \Vok & chave da Lista, Q5 (NO) \\
28 & Choque transitório de produto: efeito nulo no investimento. &
\Fno & menor, não nulo \\
29 & Queda permanente de $r$ leva o equilíbrio a um $K$ maior. &
\Vok & $q$ salta $> 1$ \\
30 & Investimento é maior quando o produto acabou de subir. &
\Vok & acelerador \\
31 & Custos de ajuste são lineares no investimento. &
\Fno & convexos ($C''>0$) \\
32 & O investimento (FBCF) é mais volátil que o consumo. &
\Vok & dados IBGE em aula \\
\bottomrule
\end{tabular}
}

\subsection*{Bloco C --- Inflação e política monetária}
{\footnotesize
\renewcommand{\arraystretch}{1.3}
\begin{tabular}{p{0.5cm} p{9.2cm} c p{3.6cm}}
\toprule
33 & $P$ sobe com $M\uparrow$, $i\uparrow$ e $Y\uparrow$. &
\Fno & $Y$ entra negativo \\
34 & $i = r + \pi^e$. & \Vok & Fisher \\
35 & Preços flexíveis: moeda não move $Y$ nem $r$. &
\Vok & neutralidade \\
36 & $\Delta M\uparrow$ permanente $\Rightarrow$ $P$ salta e $M/P$ cai. &
\Vok & Fig.\ 11.2 \\
37 & Redução permanente de $\Delta M$: $P$, $\pi^e$ e $i$ caem; $M/P$ sobe. &
\Vok & frase de ouro 5 \\
38 & Efeito liquidez: expansão monetária eleva juros curtos. &
\Fno & reduz \\
39 & Expansão monetária: curtas caem, longas sobem. &
\Vok & expectativas \\
40 & $i^n = n \cdot i^1$. & \Fno & média $+$ prêmio \\
41 & Longas mudam por expectativas de curtas futuras. &
\Vok & teoria das expectativas \\
42 & A regra de Clarida et al.\ é backward-looking. &
\Fno & forward (esperada) \\
43 & $\phi_\pi > 1$ $\Rightarrow$ o juro real sobe após choque de inflação. &
\Vok & Princípio de Taylor \\
44 & Com $\phi_\pi < 1$ a inflação converge, apenas devagar. &
\Fno & persiste e desancora \\
45 & No médio prazo, $i = r^n + \pi^*$ ($= 8\%$ na calibração). &
\Vok & desvios zeram \\
46 & Inflação alta distorce tributos de forma regressiva (Feldstein). &
\Vok & custo da inflação \\
\bottomrule
\end{tabular}
}

\subsection*{Bloco D --- Política fiscal}
{\footnotesize
\renewcommand{\arraystretch}{1.3}
\begin{tabular}{p{0.5cm} p{9.2cm} c p{3.6cm}}
\toprule
47 & Solvência: VP dos superávits primários $\geq D(0)$. &
\Vok & frase de ouro 7 \\
48 & Solvência se avalia pelo resultado corrente. &
\Fno & fluxos esperados \\
49 & Pior saldo corrente $\Rightarrow$ necessariamente mais arriscado. &
\Fno & ``necessariamente'' \\
50 & $T$ e $D$ somem da restrição das famílias; só $G$ importa. &
\Vok & ricardiana \\
51 & Corte de tributos, dado $G$: renda extra é poupada. &
\Vok & tributos futuros virão \\
52 & A equivalência vale mesmo sem o governo satisfazer sua restrição. &
\Fno & exige igualdade \\
53 & Giavazzi--Pagano: Dinamarca/Irlanda tiveram booms via expectativas. &
\Vok & contração expansionista \\
54 & Consolidação fiscal é sempre contracionista. &
\Fno & composição importa \\
\bottomrule
\end{tabular}
}

% ------------------------------------------------------------------
\section{Os 10 minutos antes da prova (última olhada)}
% ------------------------------------------------------------------

\begin{provabox}[title={Se só der para lembrar de 10 coisas}]
\begin{enumerate}[leftmargin=1.6em, itemsep=0.2em]
\item Transitório $= 1/T$; permanente $\approx 1$; amostra curta $\hat b \ll
1$, longa $\hat b \to 1$.
\item $\lambda$: significativo, $\approx 0{,}5$, por IV; rejeição
\emph{plena} exigiria $\lambda = 1$ (\Vok\ de ouro).
\item Juros $\times$ consumo: teoria \emph{positiva}, evidência \emph{fraca},
culpa do $\theta$ alto (EIS $= 1/\theta$).
\item Famílias poupam \emph{pouco} (puzzle); Carroll: desconto alto anula o
precaucional.
\item Ótimo da firma: $1 + C'(I) = q$ (\emph{iguala}); $\dot K = 0$ em
$q = 1$; $\dot q = 0$: $q = \pi(K)/r$ decrescente.
\item $K$ baixo $+$ $q$ baixo $\Rightarrow$ $q$ cai mais (\Vok\ de ouro);
choque transitório: efeito \emph{menor}, nunca nulo; juros \emph{longos}
mandam no $q$.
\item $P = M/L(i,Y)$: $Y$ entra \emph{negativo}; inflação persistente $=$
moeda; Fisher $i = r + \pi^e$.
\item Liquidez: curtas \emph{caem}, longas \emph{sobem}; $i^n =$ média das
curtas esperadas $+$ prêmio.
\item Clarida \emph{forward-looking}; $\phi_\pi > 1 \Rightarrow$ \emph{real}
sobe; $\phi_\pi < 1 \Rightarrow$ desancora; médio prazo $i = r^n + \pi^* =
8\%$.
\item Fiscal: VP dos superávits $\geq D(0)$ (esperados, não correntes!);
ricardiana: $T$/$D$ somem, \emph{G} fica; Dinamarca/Irlanda expansionistas;
dívida alta $\Rightarrow$ retorno requerido \emph{alto}.
\end{enumerate}
\end{provabox}

\vfill
{\small\color{grayrbc} Decoreba V/F --- P2 2026, Macroeconomia I (PPGEco-UFES).
Todo o conteúdo ancorado nos gabaritos oficiais (P2 2021, P2 2024), na chave
da Lista 3 2026 e nos slides das Aulas 10--13. Complementa:
\texttt{lista\_III/guia\_p2\_2026.pdf} (teoria) e
\texttt{provas/p2\_simulado\_2026.pdf} (simulado). Uso pessoal.}

\end{document}
